% Párrafos propuestos: no insertados automáticamente ni resultados nuevos.
% Mantener el idioma editorial del manuscrito mediante traducción revisada al integrar.
% Vincular por labels/secciones, no por numeración de página cambiante.

% ABSTRACT — adición de alcance; no reemplazar los resultados congelados.
La propuesta articula una hipótesis arquitectónica por capas y un programa de control interpretable. Los estudios reportados no acreditan una ventaja específica del emparejamiento polar nominado ni de la propagación inter-polar configurada. El prototipo integrado sí verifica interfaces y algunas funciones operacionales delimitadas, con utilidad práctica respaldada para ciertos componentes y no suficientemente acreditada para otros. La equivalencia de esta realización con formulaciones genéricas impide atribuir diferenciación a las coordenadas por sí solas, sin excluir futuras contribuciones de organización o aprendizaje. Estos resultados no validan la arquitectura completa ni constituyen evidencia de conciencia; esa relevancia requiere una teoría y un protocolo independientes.

% SECCIÓN 1 — contrato científico y objeto de contraste.
El objeto inmediato de evaluación es una restricción, relación o mecanismo explícito de organización, no el nombre de una capa ni la representación de dos polos. Se distinguen implementación, contribución causal, utilidad práctica y generalización; ninguna transición entre estos niveles es automática. Cada realización se vincula al mecanismo original indicando si lo realiza, aproxima, sustituye o introduce una hipótesis nueva. En particular, usar un predictor aprendido y control predictivo restringido en lugar de propagación configurada es una decisión de ingeniería, no la validación retrospectiva de las ecuaciones originales. Equilibrio dinámico significa capacidad de reorganización contextual que conserva inactividad, predominancia y coactivación, no convergencia obligatoria al punto medio ni uniformidad entre agentes.

% DISCUSIÓN/CONCLUSIÓN solicitada como §13 — insertar por heading/label verificado.
El problema principal no es únicamente la ausencia de capas, sino la falta de correspondencia suficientemente precisa entre la hipótesis polar original, sus realizaciones computacionales y las afirmaciones respaldadas por los experimentos. Los estudios descritos no acreditan una ventaja específica del emparejamiento polar nominado ni de la propagación inter-polar configurada. El prototipo integrado sí demuestra algunas funciones operacionales delimitadas, pero no establece todavía la necesidad de todos sus componentes, su generalización o el alcance completo de las capas propuestas. La equivalencia con formulaciones genéricas impide atribuir diferenciación a las coordenadas por sí solas; no excluye investigar una contribución basada en organización, aprendizaje o integración, siempre que se formule de manera falsable y se contraste con alternativas pertinentes. La siguiente etapa debe precisar esa contribución, aislar sus mecanismos y conservar separadas la evidencia de control, la evidencia arquitectónica y las hipótesis sobre conciencia.
